\section{Conclusions}
\label{sec:conclusions}

This study proposed a double-8 loading method to systematically isolate the effect of shear stress directionality on liquefaction resistance, and investigated the underlying mechanisms through DEM simulations with microscopic fabric analyses. The main findings are summarized as follows:

\begin{enumerate}
  \item The proposed loading path design eliminates the confounding factors present in conventional comparisons. The equivalent-magnitude unidirectional path ensures identical instantaneous stress magnitude to the single-8 path, isolating the effect of direction. The double-8 path further matches the rate of stress change to the single-8 path, isolating the effect of loading history (axis alternation).

  \item Liquefaction resistance follows the ordering: unidirectional $>$ single-8 $>$ double-8. Even when confounding factors are eliminated, multidirectional shear significantly reduces the number of cycles to liquefaction, confirming that shear stress directionality alone is a governing factor.

  \item The total cumulative shear work at liquefaction is approximately equal across the three loading types, indicating that directionality affects the efficiency of energy conversion to pore pressure, not the total energy threshold for liquefaction.

  \item Microscopic analyses reveal that directional changes accelerate fabric degradation: the double-8 path produces the fastest reduction in mechanical coordination number $\Zm$ and the highest development of fabric anisotropy $\Fc$, while the unidirectional path maintains the most stable contact structure.

  \item The mechanism underlying directionality-dependent liquefaction can be attributed to the disruption--realignment cycle: periodic changes in loading direction prevent the contact network from reaching a stable configuration, promoting irreversible contact loss and anisotropic fabric development.
\end{enumerate}

These results suggest that conventional unidirectional cyclic testing may underestimate the liquefaction susceptibility of soils subjected to real earthquake loading, where multidirectional shear is inherent. The proposed double-8 loading framework provides a rational basis for quantifying the additional liquefaction risk associated with shear stress directionality.
